% modules/open-targets/A2-bernstein-von-mises.tex
% -----------------------------------------------------------------------------
% Deep-dive companion to the A2 entry of the Part I agenda (modules/08-open-targets):
% asymptotic functional-inequality constants via the Bernstein--von Mises theorem.
% Compiles inside main.tex; standalone it shows cross-module \ref as "??".
% -----------------------------------------------------------------------------
\documentclass[../../main.tex]{subfiles}
\begin{document}
\section{A2 --- Asymptotic constants via Bernstein--von Mises}
\label{sec:a2}

This is the deep-dive companion to the A2 paragraph of the agenda (Section~\ref{sec:agenda}).
The target is to convert the oldest structural fact of parametric Bayesian statistics --- the
asymptotic normality of the posterior --- into sharp asymptotics for the three
functional-inequality constants. The slogan version,
\[
  \pi_n\approx N\!\big(\hat\theta_n,(nI(\theta_0))^{-1}\big)
  \quad\Longrightarrow\quad
  \CP(\pi_n)\approx\CLS(\pi_n)\approx\CTCI(\pi_n)\approx\tfrac1n\,\lmax\!\big(I(\theta_0)^{-1}\big),
\]
is correct for the Gaussian Laplace approximation itself (Family~\ref{fam:gaussian-covariance})
but is \emph{not} a corollary of the Bernstein--von Mises (BvM) theorem as usually stated. The
purpose of this section is to make the target precise, to isolate exactly where the naive
implication fails, to record that the Poincar\'e half is essentially within reach of existing
low-temperature spectral-gap theory while the log-Sobolev half is genuinely more delicate, and
to state the theorem one should actually aim to prove.

\subsection{Setup, scaling, and the precise target}
\label{subsec:a2-setup}

Fix the dimension $d$ throughout (the growing-dimension regime is deferred to
Question~\ref{q:a2-highdim}). Let $\theta_0$ be an interior point of the parameter space, the
data-generating truth, and write the (negative log-)posterior as
\[
  \pi_n(\dd\theta)\ \propto\ e^{-V_n(\theta)}\,\dd\theta,
  \qquad
  V_n(\theta)\ =\ \sum_{i=1}^n \ell_i(\theta)\ -\ \log\rho(\theta),
\]
with $\ell_i$ the per-observation negative log-likelihood and $\rho$ a prior density positive
and smooth near $\theta_0$. Let $\hat\theta_n=\arg\min_\theta V_n(\theta)$ (the MAP), let
$H_n=\Hess V_n(\hat\theta_n)$ be the posterior precision at the mode, and let $I(\theta_0)$ be
the Fisher information, so that under standard regularity $\tfrac1n H_n\to I(\theta_0)$ in
$\Prob_{\theta_0}$-probability. The classical theorem of Bernstein and von Mises
\cite{VanDerVaart1998,GhosalVanDerVaart2017} states that, after the local rescaling
\begin{equation}\label{eq:a2-rescaling}
  z\ =\ \sqrt n\,(\theta-\hat\theta_n),
  \qquad
  \tilde\pi_n\ =\ \mathrm{law\ of}\ z\ \mathrm{under}\ \pi_n,
\end{equation}
the rescaled posterior converges to a fixed Gaussian \emph{in total variation},
\begin{equation}\label{eq:a2-bvm}
  \big\|\tilde\pi_n-N\!\big(0,I(\theta_0)^{-1}\big)\big\|_{\mathrm{TV}}\ \xrightarrow{\ \Prob_{\theta_0}\ }\ 0 .
\end{equation}

The three constants scale cleanly under \eqref{eq:a2-rescaling}. If $g(z)=f(\hat\theta_n+z/\sqrt n)$
then $\Var_{\pi_n}(f)=\Var_{\tilde\pi_n}(g)$ while $\nabla_z g=n^{-1/2}\nabla_\theta f$, so testing
Definition~\ref{def:poincare} on both sides gives the exact identity
\begin{equation}\label{eq:a2-scaling}
  \CP(\pi_n)=\tfrac1n\,\CP(\tilde\pi_n),\qquad
  \CLS(\pi_n)=\tfrac1n\,\CLS(\tilde\pi_n),\qquad
  \CTCI(\pi_n)=\tfrac1n\,\CTCI(\tilde\pi_n).
\end{equation}
The $1/n$ is therefore not the content; the content is the limit of the \emph{rescaled}
constants. The target is the following.

\begin{conjecture}[Asymptotic functional-inequality constants]
\label{conj:a2}
Under regularity conditions strong enough to control the posterior tails uniformly (made
explicit in Theorem~\ref{thm:a2-target} below), for a regular parametric model in fixed
dimension $d$,
\begin{equation}\label{eq:a2-target}
  n\,\CP(\pi_n)\ \xrightarrow{\ \Prob_{\theta_0}\ }\ \lmax\!\big(I(\theta_0)^{-1}\big)
  \ =\ \frac{1}{\lmin\!\big(I(\theta_0)\big)} ,
\end{equation}
and, under additional global hypotheses (Subsection~\ref{subsec:a2-lsi}), the same limit holds
for $n\,\CLS(\pi_n)$ and $n\,\CTCI(\pi_n)$.
\end{conjecture}

Equivalently, by \eqref{eq:a2-scaling}, the claim is that the rescaled constants converge to
those of the limiting Gaussian, $\CP(\tilde\pi_n)\to\CP(N(0,I^{-1}))=\lmax(I(\theta_0)^{-1})$,
and likewise for $\CLS,\CTCI$ --- i.e.\ the BvM Gaussian computes the constants, not merely the
distribution.

\subsection{Why BvM in total variation is not enough}
\label{subsec:a2-tv}

The constants are far more sensitive than total-variation distance: they see tails, remote
mass, gradients, and global geometry, none of which TV controls. Convergence
\eqref{eq:a2-bvm} does not by itself control even posterior covariances, let alone
functional-inequality constants \cite{KasprzakGiordanoBroderick2025}.

\begin{warning}[A TV-BvM sequence with diverging constants]
\label{warn:a2-tv-fails}
Consider the contamination sequence on $\R$
\[
  \mu_n\ =\ (1-\eps_n)\,N(0,1)\ +\ \eps_n\,N(a_n,1),
  \qquad \eps_n\to0,\quad \eps_n a_n^2\to\infty .
\]
Then $\|\mu_n-N(0,1)\|_{\mathrm{TV}}\le\eps_n\to0$, so $\mu_n$ satisfies a BvM-type limit in TV.
Yet its variance is $\Var_{\mu_n}(x)=1+\eps_n a_n^2-\eps_n^2a_n^2\to\infty$, and the linear test
$f(x)=x$ gives the lower bound $\CP(\mu_n)\ge\Var_{\mu_n}(x)\to\infty$
(Remark~\ref{rem:rayleigh-lower-only}). The same remote sliver of mass inflates $\CLS$ and
$\CTCI$. Hence \eqref{eq:a2-bvm} alone is consistent with all three constants diverging.
\end{warning}

The lesson is structural: to extract constants one needs a \emph{strong} Laplace/BvM statement
that adds (i) a uniform local quadratic approximation of $V_n$ near $\hat\theta_n$ and (ii)
tail / no-bottleneck control ruling out remote mass of the kind in
Warning~\ref{warn:a2-tv-fails}. The clean device for (ii) is bounded perturbation: if the
rescaled posterior is a globally $o_P(1)$-oscillation reweighting of its BvM Gaussian, then
Holley--Stroock (Theorem~\ref{thm:holley-stroock}) transfers the Gaussian constants with a
multiplicative error $e^{\osc}=1+o_P(1)$, for $\CP$ and $\CLS$ simultaneously. The whole
difficulty of A2 is to discharge (i)--(ii) under hypotheses weak enough to cover real models
yet strong enough to defeat Warning~\ref{warn:a2-tv-fails}.

\subsection{The Poincar\'e half is within reach: low-temperature spectral gaps}
\label{subsec:a2-poincare}

Write the posterior as a low-temperature Gibbs measure. With $t=1/n$ and
$f_n(\theta)=\tfrac1n V_n(\theta)\to\ell(\theta)$ (the population risk, minimized at $\theta_0$
with $\Hess\ell(\theta_0)=I(\theta_0)$),
\[
  \pi_n(\dd\theta)\ \propto\ \exp\!\big(-f_n(\theta)/t\big)\,\dd\theta,
  \qquad t=1/n\to0 .
\]
This is exactly the object of low-temperature functional-inequality theory. Chewi and Stromme
\cite{ChewiStromme2024Ballistic} prove that for $\mu_t\propto e^{-f/t}$ with a unique
nondegenerate minimizer $x_\star$,
\begin{equation}\label{eq:a2-cs-poincare}
  \lim_{t\to0}\frac{\CP(\mu_t)}{t}\ =\ \frac{1}{\lmin\!\big(\Hess f(x_\star)\big)} ,
\end{equation}
the inverse smallest Hessian eigenvalue at the minimizer --- a purely \emph{local} quantity.
Substituting $f=\ell$, $x_\star=\theta_0$, $t=1/n$ reproduces the right-hand side of
\eqref{eq:a2-target} term for term. Thus the Poincar\'e statement is not a fantasy: it is the
statistical shadow of \eqref{eq:a2-cs-poincare}.

\begin{remark}[What transfer to the random posterior still requires]
\label{rem:a2-transfer}
Equation~\eqref{eq:a2-cs-poincare} is a deterministic statement about a fixed potential $f$.
The posterior potential $f_n=\tfrac1nV_n$ is \emph{random} (empirical-likelihood fluctuations),
its minimizer $\hat\theta_n$ is random and only converges to $\theta_0$, and the prior
contributes an $O(1/n)$ perturbation. Turning \eqref{eq:a2-cs-poincare} into
\eqref{eq:a2-target} therefore needs the low-temperature theorem in a form uniform over a
shrinking neighbourhood of potentials around $\ell$, plus a global tail bound excluding remote
posterior mass (the Warning~\ref{warn:a2-tv-fails} pathology). Related low-temperature toolkits
that could supply the global control are Lyapunov-potential criteria
\cite{ChenSridharan2024Lyapunov}, local log-Polyak--{\L}ojasiewicz Poincar\'e bounds
\cite{GongHeShen2025LogPL}, and the constrained-Gibbs posterior Poincar\'e inequalities of
\cite{ConstrainedGibbs2026}, which already prove $O(1/n)$ local Poincar\'e constants on a
high-probability good set for logistic regression, Poisson linear models, and Gaussian mixtures.
\end{remark}

\subsection{The log-Sobolev half is genuinely more delicate}
\label{subsec:a2-lsi}

This is where the naive paragraph is most misleading. For a Gaussian $\CLS=\CP$
(Family~\ref{fam:gaussian-covariance}), so the local-Hessian answer is correct in the
Gaussian-perturbative regime. But for a non-Gaussian low-temperature Gibbs measure the
log-Sobolev constant depends on the \emph{global} landscape, not only on the curvature at the
minimizer. The companion theorem of \cite{ChewiStromme2024Ballistic} identifies the ballistic
limit of the log-Sobolev constant with the \emph{Polyak--{\L}ojasiewicz} (PL) constant of $f$:
\begin{equation}\label{eq:a2-cs-lsi}
  \lim_{t\to0}\frac{\CLS(\mu_t)}{t}\ =\ \frac{1}{\mu_{\mathrm{PL}}(f)},
  \qquad
  \tfrac12\norm{\nabla f(x)}^2\ \ge\ \mu_{\mathrm{PL}}(f)\,\big(f(x)-\min f\big)\ \ \forall x,
\end{equation}
where $\mu_{\mathrm{PL}}(f)$ is the largest valid PL constant. Since
$\mu_{\mathrm{PL}}(f)\le\lmin(\Hess f(x_\star))$ in general, with equality only when the global
landscape is as favourable as the local curvature, one has
\[
  \lim_{t\to0}\frac{\CLS(\mu_t)}{t}\ \ge\ \lim_{t\to0}\frac{\CP(\mu_t)}{t},
\]
consistent with $\CP\le\CLS$ (Theorem~\ref{thm:lsi-implies-poincare}), and the inequality is
\emph{strict} whenever the landscape away from $\theta_0$ is shallow relative to $I(\theta_0)$.

\begin{warning}[The LSI limit need not be Fisher-local]
\label{warn:a2-lsi}
Consequently the clean statement
\(
  n\,\CLS(\pi_n)\to\lmax(I(\theta_0)^{-1})
\)
is valid only in special regimes --- when the posterior is a global $o_P(1)$-oscillation
perturbation of its BvM Gaussian (Holley--Stroock with vanishing oscillation), or when the
negative log-likelihood is globally strongly convex / PL with constant matching $I(\theta_0)$.
In general the correct asymptotic is
\(
  n\,\CLS(\pi_n)\to 1/\mu_{\mathrm{PL}}(\ell),
\)
and $1/\mu_{\mathrm{PL}}(\ell)$ need not equal $\lmax(I(\theta_0)^{-1})$. Multiple or shallow
remote near-minimizers change the LSI behaviour even when the BvM Gaussian is unchanged
\cite{BenNejma2026Lojasiewicz}.
\end{warning}

\subsection{Transport follows from a sharp LSI, not from BvM}
\label{subsec:a2-transport}

The Talagrand $T_2$ constant inherits both the good and the bad news. Whenever the
log-Sobolev comparison is sharp, $n\,\CLS(\pi_n)\to\lmax(I^{-1})$, Otto--Villani
(Theorem~\ref{thm:otto-villani}) gives the matching upper bound
$n\,\CTCI(\pi_n)\le\lmax(I^{-1})(1+o_P(1))$, while the linear-test / covariance lower bound
(Remark~\ref{rem:rayleigh-lower-only}) supplies the reverse inequality under enough moment
control, so $n\,\CTCI(\pi_n)\to\lmax(I(\theta_0)^{-1})$. Outside the Gaussian-perturbative
regime the $T_2$ asymptotic degrades exactly as the LSI one does (Warning~\ref{warn:a2-lsi}).
In all cases this is a consequence of a \emph{strong} LSI comparison, never of TV-BvM alone.

\subsection{The theorem to aim for}
\label{subsec:a2-theorem}

The honest statement is not ``BvM implies the constants'' but a strong-Laplace statement that
adds the missing tail and perturbation control.

\begin{theorem}[Target shape]
\label{thm:a2-target}
Let $d$ be fixed, $\theta_0$ interior, the prior density positive and smooth near $\theta_0$,
and the model regular with $\tfrac1n H_n\to I(\theta_0)\succ0$ in $\Prob_{\theta_0}$-probability.
Suppose that after the whitening $z=H_n^{1/2}(\theta-\hat\theta_n)$ the rescaled posterior
$\tilde\pi_n$ admits the representation
\[
  \frac{\dd\tilde\pi_n}{\dd\gamma}(z)\ =\ \frac{e^{-r_n(z)}}{\int e^{-r_n}\dd\gamma},
  \qquad \gamma=N(0,\Id_d),
\]
with a global oscillation bound $\osc(r_n)=o_P(1)$ (equivalently, any tail/no-bottleneck
condition strong enough to defeat Warning~\ref{warn:a2-tv-fails} and supply
\eqref{eq:a2-cs-poincare} uniformly). Then
\[
  \CP(\pi_n)\ =\ (1+o_P(1))\,\lmax(H_n^{-1})\ =\ \frac{1+o_P(1)}{n}\,\lmax\!\big(I(\theta_0)^{-1}\big),
\]
and, under the same hypothesis (which makes the perturbation Gaussian-comparable for LSI too),
\[
  \CLS(\pi_n),\ \CTCI(\pi_n)\ =\ \frac{1+o_P(1)}{n}\,\lmax\!\big(I(\theta_0)^{-1}\big).
\]
\end{theorem}

\noindent The proof skeleton is: whitening reduces the BvM Gaussian to the standard $\gamma$,
whose constants are $1$; Holley--Stroock (Theorem~\ref{thm:holley-stroock}) transfers them to
$\tilde\pi_n$ with factor $e^{\osc(r_n)}=1+o_P(1)$ for $\CP$ and $\CLS$; Otto--Villani
(Theorem~\ref{thm:otto-villani}) and the covariance lower bound
(Remark~\ref{rem:rayleigh-lower-only}) pin $\CTCI$; the scaling identity \eqref{eq:a2-scaling}
and $\tfrac1nH_n\to I(\theta_0)$ restore the $1/n$ and the Fisher information. The mathematical
weight is entirely in verifying the oscillation/tail hypothesis for concrete model classes ---
which is where the perturbative ``$\osc=o_P(1)$'' assumption is too strong for the Poincar\'e
conclusion (there \eqref{eq:a2-cs-poincare} only needs \emph{local} uniformity plus a tail
bound) and is genuinely needed for the log-Sobolev conclusion.

\subsection{Sub-targets and what remains open}
\label{subsec:a2-open}

The single conjecture splits into four research questions of increasing difficulty, the first
of which is essentially a transfer exercise and the last of which is wide open.

\begin{question}[The weakest hypothesis for the Poincar\'e limit]
\label{q:a2-poincare}
Identify the weakest posterior assumptions under which $n\,\CP(\pi_n)\to\lmax(I(\theta_0)^{-1})$.
The deterministic low-temperature analogue \eqref{eq:a2-cs-poincare} is settled
\cite{ChewiStromme2024Ballistic}; what is missing is the random-posterior theorem with
empirical-likelihood fluctuations, a prior, and the common statistical models, with a clean
tail condition replacing the strong $\osc=o_P(1)$ assumption of Theorem~\ref{thm:a2-target}.
This is the entry point and should be provable with current tools
\cite{ConstrainedGibbs2026,GongHeShen2025LogPL,ChenSridharan2024Lyapunov}.
\end{question}

\begin{question}[When is the LSI limit Fisher-local?]
\label{q:a2-lsi}
Characterize the models for which $n\,\CLS(\pi_n)\to\lmax(I(\theta_0)^{-1})$ (Fisher-local)
versus those for which the limit is the global-landscape quantity
$1/\mu_{\mathrm{PL}}(\ell)$ (Warning~\ref{warn:a2-lsi}). Give checkable sufficient conditions
--- e.g.\ global log-concavity of the population risk, or a quantitative
single-well/PL hypothesis --- under which the two coincide, and exhibit a natural statistical
model where they provably differ \cite{BenNejma2026Lojasiewicz}.
\end{question}

\begin{question}[Multimodal and weakly identified posteriors]
\label{q:a2-multimodal}
For mixtures, label switching, weak identification, and phase transitions the posterior is
multimodal and $\CP(\pi_n)$ can be \emph{exponentially} large in $n$ despite local Gaussianity
around each mode --- Eyring--Kramers / metastability behaviour rather than \eqref{eq:a2-target}.
Determine when the blow-up is physical and when it is a symmetry artifact removable by passing
to the quotient (the symmetry-quotient escape of Remark~\ref{rem:label-switching}, the
$e^{c\Delta^2}$ barrier of Theorem~\ref{thm:separated-mixture-lower}). This connects A2 to A5
(Section~\ref{sec:a5}, Conjecture~\ref{conj:a5-metastable}), where the same blow-up is shown to
dissolve on the symmetry quotient.
\end{question}

\begin{question}[Growing dimension $d=d_n$]
\label{q:a2-highdim}
Extend the sharp constants to the regime $d=d_n\to\infty$. MCMC mixing analyses already track
the joint $(d,n,\kappa)$ dependence using large-sample posterior Gaussianity
\cite{TangYang2024MALA}, and polynomial-time Langevin guarantees for high-dimensional
posteriors are known but deliberately avoid full asymptotic Gaussian approximation, relying
instead on local curvature and global log-concave comparison \cite{NicklWang2024Langevin};
exact Fisher-information asymptotics for $\CP,\CLS,\CTCI$ in growing dimension are far less
settled. Log-concave BvM theory \cite{Brunel2026LogConcaveBvM} and the finite-sample Laplace
bounds of \cite{KasprzakGiordanoBroderick2025} are the natural starting points.
\end{question}

\begin{remark}[Relation to the rest of the agenda and to numerics]
\label{rem:a2-relations}
A2 is the asymptotic counterpart of the finite-sample flagship A1
(Section~\ref{sec:a1}, Question~\ref{q:glm-data-informed}, Conjecture~\ref{conj:a1}): A1 seeks a
data-informed bound at fixed $n$ via Brascamp--Lieb (Theorem~\ref{thm:brascamp-lieb}) and
posterior concentration, while A2 seeks the exact $n\to\infty$ limit; a sharp A1 bound that is
tight in the bulk should reproduce \eqref{eq:a2-target} on letting $n\to\infty$, and consistency
between the two is a useful internal check. The three remaining siblings are the regimes in which
the clean limit \eqref{eq:a2-target} is qualified: the BvM Gaussian comparison presupposes light
posterior tails and breaks for the heavy-tailed posteriors of A3 (Section~\ref{sec:a3},
Warning~\ref{warn:a3-likelihood}), where weighted rather than classical constants govern the
limit; the asymptotic scale $\tfrac1n\lmax(I(\theta_0)^{-1})$ is exactly the posterior-scale
target consumed by the variational-inference certificates of A4 (Section~\ref{sec:a4},
Conjecture~\ref{conj:a4}); and the multimodal blow-up of Question~\ref{q:a2-multimodal} is
resolved on the symmetry quotient of A5 (Section~\ref{sec:a5}, Conjecture~\ref{conj:a5-metastable}),
where the BvM asymptotic re-emerges. Empirically, the lower bound $\CP(\pi_n)\ge\lmax(\Cov_{\pi_n})$
(Remark~\ref{rem:rayleigh-lower-only}) and statistical estimators of the Poincar\'e constant
from samples \cite{PillaudVivienEtAl2020} give a numerical handle on the conjectured limit,
with the ratio of a curvature/perturbation upper bound to $\lmax(\Cov_{\pi_n})$ serving as the
sharpness diagnostic.
\end{remark}
\end{document}
